% =====================================================================
%  MỤC 6 — KẾT LUẬN
% =====================================================================
\section{Kết luận}

Từ kết quả phân tích chuỗi số liệu mưa giờ tại trạm \TenTram{} giai đoạn \GiaiDoan{} bằng các công cụ Python (pandas, numpy, scipy, matplotlib), có thể rút ra các kết luận sau:

\begin{enumerate}[label=(\arabic*), leftmargin=1.2cm]
  \item \textbf{Thống kê mưa ngày, tháng, năm và mùa (câu 1).} \TenTram{} là khu vực mưa rất lớn với lượng mưa năm trung bình 4\,192,5 mm (2\,425,5--5\,567,9 mm, CV 23,3\%). Mùa mưa IX--XII đóng góp trung bình 68\% lượng mưa năm; tháng XI mưa nhiều nhất (trung bình 1\,167,0 mm), tháng II và IV ít mưa nhất (khoảng 93 mm). Mưa ngày phân bố lệch phải mạnh (trung bình 12,39 mm, CV 278\%, độ lệch 8,22); tần suất ngày mưa $\geq 1$ mm là 46,6\%, các ngày mưa trên 100 mm chỉ chiếm 2\% số ngày nhưng đóng góp 31,8\% tổng lượng mưa. Mưa giờ mùa khô mang tính đối lưu với cực đại buổi chiều, còn mưa mùa mưa kéo dài đều trong ngày.
  \item \textbf{Mưa ngày lớn nhất và mưa 5 ngày lớn nhất năm (câu 2).} Rx1day trung bình 364,8 mm, lớn nhất 781,6 mm (05/11/2017); Rx5day trung bình 677,7 mm, lớn nhất 1\,255,4 mm (03--07/11/2017). Mọi Rx1day đều xảy ra trong các tháng IX--XI. Rx5day thường gấp khoảng 2 lần Rx1day. Ngày Rx1day trùng ngày trượt lở ở 8/10 năm, giai đoạn Rx5day chứa ngày trượt lở ở 9/10 năm.
  \item \textbf{Ngưỡng mưa rất lớn và cực kỳ lớn (câu 3).} Với mưa giờ, R95p = 20,1 mm/giờ và R99p = 39,4 mm/giờ; với mưa ngày, R95p = 89,2 mm/ngày và R99p = 209,0 mm/ngày (mùa mưa: 154,6 và 338,9 mm/ngày). 5\% số ngày mưa vượt R95p đóng góp 34,3\% tổng lượng mưa ngày mưa. 27/28 ngày có trượt lở có lượng mưa vượt R95p ngày, nên ngưỡng này có ý nghĩa thực tiễn cho cảnh báo trượt lở.
  \item \textbf{Mưa cực đại theo chu kỳ lặp lại (câu 4).} Phân bố Gumbel được chọn cho cả hai chuỗi. Mưa giờ cực đại ứng với chu kỳ 10, 20, 50, 100 năm lần lượt là 85,0; 94,0; 105,8 và 114,5 mm; mưa ngày cực đại lần lượt là 564,3; 651,7; 764,9 và 849,8 mm. Mức 10--20 năm nhất quán giữa các phương pháp (sai khác dưới 10\%); mức 50--100 năm có độ bất định lớn (khoảng tin cậy 95\% của mưa ngày 100 năm là 548,5--1\,141,1 mm) và nhạy với giá trị cực trị năm 2017.
  \item \textbf{Sản phẩm trực quan hóa.} Toàn bộ kết quả được công bố thành website tra cứu công khai \WebURL{} cùng mã nguồn mở tại \GitHubURL. Người dùng có thể tra cứu lượng mưa của bất kỳ ngày, giờ nào trong 10 năm, xem lại các bảng và biểu đồ của 4 câu hỏi dưới dạng tương tác, quy đổi giữa chu kỳ lặp lại và lượng mưa, và tải dữ liệu về. Quy trình kiểm tra tự động bảo đảm số liệu trên website luôn khớp với báo cáo.
\end{enumerate}

\paragraph{Hạn chế và hướng phát triển.} Chuỗi số liệu chỉ dài 10 năm, của một trạm, và thiếu toàn bộ tháng 11/2009 (tháng mưa lớn nhất) cùng phần lớn tháng 01/2017 và các ngày mùng 1 hằng tháng; vì vậy các ước lượng cho chu kỳ 50--100 năm chỉ mang tính tham khảo. Hướng phát triển gồm: bổ sung chuỗi số liệu dài hơn và các trạm lân cận để phân tích tần suất vùng; phân tích mưa thời đoạn ngắn (1, 3, 6, 12 giờ) để xây dựng đường cong cường độ -- thời gian -- tần suất (IDF); xây dựng ngưỡng mưa kích hoạt trượt lở kết hợp lượng mưa tích lũy nhiều ngày; và mở rộng website cho nhiều trạm đo, cập nhật số liệu mới định kỳ để phục vụ tra cứu, cảnh báo mưa lớn tại địa phương.
